% GENERATED FILE. Edit canonical lessons, then run npm run generate:books.
% canonical-source-hash: fnv1a64:7351fee2e5dd6a41
% canonical-lessons: ZH-C459-gudan, ZH-C459-shiqing, ZH-W459-qing-feeling, ZH-C459-ganqing, ZH-C459-xinlihua, ZH-C459-weiqu

\chapter{On My Mind}
\label{ch:zh-on-my-mind}
\hlchaptermodality{459}

\begin{chapteropening}
\textbf{By the end of this chapter:} \emph{I can say I feel lonely, that I have a lot to do, talk about feelings between people and what is on my mind.}

Complete the last lesson of chapter 459: I can say I feel lonely, that I have a lot to do, talk about feelings between people and what is on my mind.
\end{chapteropening}

\section[gūdān]{gūdān — lonely; on one's own}

\label{lesson:ZH-C459-gudan}



\begin{quote}
Before the new one: say the Mandarin for relaxed and at ease, then the Mandarin for fed up.
\end{quote}

\subsection*{You'll want to know: gūdān}
\emph{gūdān} — \textquotedblleft{}lonely; on one's own\textquotedblright{}.

\begin{quote}
\emph{wǒ yí ge rén zhù le liǎng nián yǒudiǎnr gūdān} — I've lived on my own for two years, and I'm a bit lonely.
\end{quote}

Its characters come later; for now it is a word for the ear.

\begin{cousinweb}
\emph{gūdān} is built from \emph{gū} 'an orphan' and \emph{dān} 'single', both inherited from Middle Chinese, \emph{ku} and \emph{tan} (Guangyun, in Baxter's transcription), which glosses the first 'a fatherless child' and the second as single against double.
\end{cousinweb}

\begin{grammarlens}[title={What you've built}]
One more word for saying how it was.
\end{grammarlens}

\subsection*{Your turn}
\begin{itemize}
\raggedright
  \item \emph{Say it:} \emph{gūdān}
  \item \emph{Say it:} \emph{gūdān}, once more
  \item \emph{Say it:} \emph{wǒ yí ge rén zhù le liǎng nián yǒudiǎnr gūdān} — I've lived on my own for two years, and I'm a bit lonely.
  \item \emph{Recall:} say the Mandarin for relaxed and at ease, then the Mandarin for fed up, then say \emph{gūdān} again
\end{itemize}

\begin{tcolorbox}[breakable,colback=teal!4,colframe=teal!35!black,title={Before you move on}]
What does gūdān mean? (\textquotedblleft{}Lonely; on one's own\textquotedblright{}.) Say it once more.
\end{tcolorbox}
\section[\texorpdfstring{\zh{事情} (shìqing)}{事情 (shìqing)}]{\zh{事情} (shìqing) — matter; affair}

\label{lesson:ZH-C459-shiqing}



\begin{quote}
Before the new one: say the Mandarin for fed up, then the Mandarin for lonely.
\end{quote}

\subsection*{You'll want to know: \zh{事情}}
\textbf{\zh{事情}} — \emph{shìqing} — \textquotedblleft{}matter; affair\textquotedblright{}.

\begin{quote}
\textbf{\zh{今天我有很多事情要做}} \emph{jīntiān wǒ yǒu hěn duō shìqing yào zuò} — I've got a lot of things to do today.
\end{quote}

The new character, in pieces: \zh{情} — eleven strokes in three pieces: on the left the heart piece, a dot, a second dot and a long vertical; at the top right a bar, a second bar, a vertical through them and a long bar; and under them \zh{月} \emph{yuè}, four strokes. The next lesson puts it in your hand.

\begin{cousinweb}
\textbf{\zh{事情}} (\emph{shìqing}) is built from \zh{事} 'an affair' and \zh{情} 'feeling', both inherited from Middle Chinese, \emph{dzriH} and \emph{dzjeng} (Guangyun, in Baxter's transcription). The dictionary glosses the first 'to employ' and 'to set up', and explains the second, after an older book, as what a person desires; in this word the second is said light, and why it joins the first is not settled.
\end{cousinweb}

\begin{grammarlens}[title={What you've built}]
One more word for saying how it was.
\end{grammarlens}

\subsection*{Your turn}
\begin{itemize}
\raggedright
  \item \emph{Say it:} \emph{shìqing}
  \item \emph{Say it:} \emph{shìqing}, once more
  \item \emph{Say it:} \emph{jīntiān wǒ yǒu hěn duō shìqing yào zuò} — I've got a lot of things to do today.
  \item \emph{Recall:} say the Mandarin for fed up, then the Mandarin for lonely, then say \emph{shìqing} again
\end{itemize}

\begin{tcolorbox}[breakable,colback=teal!4,colframe=teal!35!black,title={Before you move on}]
What does \zh{事情} mean? (\textquotedblleft{}Matter; affair\textquotedblright{}.) Say it once more.
\end{tcolorbox}
\section[\texorpdfstring{\zh{情} (qíng)}{情 (qíng)}]{\zh{情} — the heart piece, then a top over \zh{月}}

\label{lesson:ZH-W459-qing-feeling}



\begin{quote}
The word first, then the character inside it.

\begin{itemize}
\raggedright
  \item \emph{Recall:} say \emph{shìqing} — \textbf{R1}, one lesson back, before it decays
\end{itemize}

\begin{quote}
\zh{月} \zh{事}
\end{quote}

\emph{Read it:} each one aloud, with nothing said first — \emph{yuè}, \emph{shì}
\end{quote}

\subsection*{Script — the heart piece, then a top over \zh{月}}
\hlblockfigure{\detokenize{figures/ZH-W459-qing-feeling-filmstrip.pdf}}{How \zh{情} is written, stroke by stroke}

\begin{quote}
\zh{情}
\end{quote}

\textbf{Three pieces: the left first, then the top right, then the bottom right.} Count the pieces, not the strokes.

\textbf{The left piece, three strokes: the heart piece.} A short dot on the left, a short dot on the right falling to the right, then a long vertical down between them from top to foot.

\textbf{The top right, four strokes.} A bar from left to right, a second bar under it, a vertical down through both, then a long bar under them, wider than the two above.

\textbf{The bottom right, \zh{月}, four strokes.} The stroke down the left that curves to the left at the foot, the stroke that runs right, turns down and hooks, then the two short bars inside.

3 + 4 + 4 = eleven.

\subsection*{Your turn}
\begin{itemize}
\raggedright
  \item \emph{Write it:} \zh{情}, all eleven, with the model in front of you
  \item \emph{Cover:} write it again from memory, and count the pieces, not the strokes
\end{itemize}

\begin{tcolorbox}[breakable,colback=teal!4,colframe=teal!35!black,title={Before you move on}]
How many pieces does \zh{情} have? (\textbf{Three: three strokes, four and four}.) Which piece comes last? (\textbf{\zh{月}, at the bottom right}.)

Source: \href{https://raw.githubusercontent.com/chanind/hanzi-writer-data/68d10a4b21150cae5e1ebbd223eed289cf32d90c/data/\%E6\%83\%85.json}{Hanzi Writer Data \zh{情}, pinned stroke paths and medians}.
\end{tcolorbox}
\section[\texorpdfstring{\zh{感情} (gǎnqíng)}{感情 (gǎnqíng)}]{\zh{感情} (gǎnqíng) — affection between people; feelings for each other}

\label{lesson:ZH-C459-ganqing}



\begin{quote}
Before the new one: say the Mandarin for lonely, then the Mandarin for a matter to deal with.
\end{quote}

\subsection*{You'll want to know: \zh{感情}}
\textbf{\zh{感情}} — \emph{gǎnqíng} — \textquotedblleft{}affection between people; feelings for each other\textquotedblright{}.

\begin{quote}
\emph{wǒmen rènshi le shí nián gǎnqíng hěn hǎo} — We've known each other for ten years, and we're very close.
\end{quote}

\begin{cousinweb}
\textbf{\zh{感情}} (\emph{gǎnqíng}) is built from \zh{感} 'to move' and \zh{情} 'feeling', both inherited from Middle Chinese, \emph{komX} and \emph{dzjeng} (Guangyun, in Baxter's transcription), which glosses the first 'to move'.
\end{cousinweb}

\begin{grammarlens}[title={What you've built}]
One more word for saying how it was.
\end{grammarlens}

\subsection*{Your turn}
\begin{itemize}
\raggedright
  \item \emph{Say it:} \emph{gǎnqíng}
  \item \emph{Say it:} \emph{gǎnqíng}, once more
  \item \emph{Say it:} \emph{wǒmen rènshi le shí nián gǎnqíng hěn hǎo} — We've known each other for ten years, and we're very close.
  \item \emph{Recall:} say the Mandarin for lonely, then the Mandarin for a matter to deal with, then say \emph{gǎnqíng} again
\end{itemize}

\begin{tcolorbox}[breakable,colback=teal!4,colframe=teal!35!black,title={Before you move on}]
What does \zh{感情} mean? (\textquotedblleft{}Affection between people; feelings for each other\textquotedblright{}.) Say it once more.
\end{tcolorbox}
\section[\texorpdfstring{\zh{心里话} (xīnlihuà)}{心里话 (xīnlihuà)}]{\zh{心里话} (xīnlihuà) — what is really on one's mind; words from the heart}

\label{lesson:ZH-C459-xinlihua}



\begin{quote}
Before the new one: say the Mandarin for a matter to deal with, then the Mandarin for affection between people.
\end{quote}

\subsection*{You'll want to know: \zh{心里话}}
\textbf{\zh{心里话}} — \emph{xīnlihuà} — \textquotedblleft{}what is really on one's mind; words from the heart\textquotedblright{}.

\begin{quote}
\emph{zhè shì wǒ de xīnlihuà bàituō nǐ bié gàosu biérén} — This is what I really think, so please don't tell anyone.
\end{quote}

\begin{cousinweb}
\textbf{\zh{心里话}} (\emph{xīnlihuà}) is built from \zh{心} 'heart', \zh{里} 'inside' and \zh{话} 'words', all inherited from Middle Chinese, \emph{sim}, \emph{liX} and \emph{hwaejH} (Guangyun, in Baxter's transcription), which glosses the second 'inside' and the third 'speech'. The middle syllable is said light.
\end{cousinweb}

\begin{grammarlens}[title={What you've built}]
One more word for saying how it was.
\end{grammarlens}

\subsection*{Your turn}
\begin{itemize}
\raggedright
  \item \emph{Say it:} \emph{xīnlihuà}
  \item \emph{Say it:} \emph{xīnlihuà}, once more
  \item \emph{Say it:} \emph{zhè shì wǒ de xīnlihuà bàituō nǐ bié gàosu biérén} — This is what I really think, so please don't tell anyone.
  \item \emph{Recall:} say the Mandarin for a matter to deal with, then the Mandarin for affection between people, then say \emph{xīnlihuà} again
\end{itemize}

\begin{tcolorbox}[breakable,colback=teal!4,colframe=teal!35!black,title={Before you move on}]
What does \zh{心里话} mean? (\textquotedblleft{}What is really on one's mind; words from the heart\textquotedblright{}.) Say it once more.
\end{tcolorbox}
\section[wěiqu]{wěiqu — wronged; hurt by unfair blame}

\label{lesson:ZH-C459-weiqu}



\begin{quote}
Before the new one: say the Mandarin for affection between people, then the Mandarin for what is really on one's mind.
\end{quote}

\subsection*{You'll want to know: wěiqu}
\emph{wěiqu} — \textquotedblleft{}wronged; hurt by unfair blame\textquotedblright{}.

\begin{quote}
\emph{biérén bù xiāngxìn wǒ wǒ gǎnjué hěn wěiqu} — Nobody believes me, and I feel really hard done by.
\end{quote}

Its characters come later; for now it is a word for the ear.

\begin{cousinweb}
\emph{wěiqu} is built from \emph{wěi} and \emph{qū}, both inherited from Middle Chinese, \emph{'jweX} and \emph{khjut} (Guangyun, in Baxter's transcription), which glosses both as bent; the sense of a person feeling wronged is not given there.
\end{cousinweb}

\begin{grammarlens}[title={What you've built}]
One more word for saying how it was.
\end{grammarlens}

\subsection*{Your turn}
\begin{itemize}
\raggedright
  \item \emph{Say it:} \emph{wěiqu}
  \item \emph{Say it:} \emph{wěiqu}, once more
  \item \emph{Say it:} \emph{biérén bù xiāngxìn wǒ wǒ gǎnjué hěn wěiqu} — Nobody believes me, and I feel really hard done by.
  \item \emph{Recall:} say the Mandarin for affection between people, then the Mandarin for what is really on one's mind, then say \emph{wěiqu} again
\end{itemize}

\begin{tcolorbox}[breakable,colback=teal!4,colframe=teal!35!black,title={Before you move on}]
What does wěiqu mean? (\textquotedblleft{}Wronged; hurt by unfair blame\textquotedblright{}.) Say it once more.
\end{tcolorbox}
